\PassOptionsToClass{onecolumn}{extarticle}
\documentclass[9pt,a4paper,twoside]{../rho-class/rho}
\usepackage{ctex}
\usepackage{amsmath}
\usepackage{amssymb}
\usepackage{siunitx}
\usepackage{hyperref}
\usepackage{indentfirst}

\setbool{rho-abstract}{true} % Set false to hide the abstract
\setbool{corres-info}{false} % Set false to hide the corresponding author section
\setbool{linenumbers}{false} % Set false to hide the line numbering

\sisetup{
    mode = math, 
    reset-math-version = false, 
}
\setlength{\parindent}{2em}
\linespread{1.2}

%----------------------------------------------------------
% TITLE & AUTHORS
%----------------------------------------------------------

\journalname{大学物理-基础实验 B 实验设计}
\title{单摆法测重力加速度 \  实验设计}

\author{崔俊熙 PB25000124}

\dates{2026-03-24}
\leadauthor{崔俊熙 PB25000124}
\footinfo{大学物理-基础实验 B 实验设计}
\smalltitle{单摆法测重力加速度 \  实验设计}
\theday{2026-03-24}

% ----------------------------------------------------------
% ABSTRACT
% ----------------------------------------------------------

\begin{abstract}
    本实验方案设计中,根据单摆周期公式,写出了单摆法测量重力加速度的不确定度合成公式;求出了需要连续测量的周期数,使得不确定度优于$\qty{1}{\percent}$.
\end{abstract}

\keywords{单摆;重力加速度;不确定度分析;周期测量}

% ----------------------------------------------------------
% DOCUMENT CONTENT
% ----------------------------------------------------------

\begin{document}

\onecolumn
\maketitle
\thispagestyle{firststyle}

\section{不确定度分析}

\subsection{运动学原理}
在摆角$\theta < 5^\circ$,忽略空气浮力及摆线质量的一级近似条件下,单摆运动可视为简谐运动,其周期公式为
\begin{equation}
    T = 2\pi\sqrt{\frac{l}{g}}
\end{equation}

其中,$l$为有效摆长(摆线长与摆球半径之和),$T$为摆动周期.

由此可导出重力加速度计算公式
\begin{equation}
    g = \frac{4\pi^2 l}{T^2}
\end{equation}

\subsection{不确定度分析推导}
对重力加速度公式进行偏导分析
\begin{equation}
    \frac{\partial g}{\partial l} = \frac{4\pi^2}{T^2} = \frac{g}{l}, \quad \frac{\partial g}{\partial T} = -\frac{8\pi^2 l}{T^3} = -\frac{2g}{T}
\end{equation}

根据误差传递公式,合成不确定度$u_g$为
\begin{equation}
    u_g = \sqrt{\left( \frac{\partial g}{\partial l} u_l \right)^2 + \left( \frac{\partial g}{\partial T} u_T \right)^2} = g \sqrt{\left( \frac{u_l}{l} \right)^2 + \left( \frac{2u_T}{T} \right)^2}
\end{equation}

即相对不确定度传递公式为
\begin{equation}
    \frac{u_g}{g} = \sqrt{\left( \frac{u_l}{l} \right)^2 + \left( \frac{2u_T}{T} \right)^2}
\end{equation}

\begin{rhoenv}[frametitle=注意.]
    \begin{itemize}
        \item 有效摆长的不确定度来自钢卷尺测量摆线的误差,以及游标卡尺测量小球直径的误差.
              其中钢卷尺的误差为$\qty{0.2}{\centi\metre}$,游标卡尺的误差为$\qty{0.02}{\centi\metre}$.
              因此$u_l \approx \sqrt{(\qty{0.2}{\centi\metre})^2 + (\qty{0.02}{\centi\metre})^2} \approx \qty{0.2}{\centi\metre}$.

        \item 周期的不确定度来自秒表的误差以及实验人员的反应时间.
              其中秒表的误差为$\qty{0.1}{\second}$,实验人员的反应时间误差为$\qty{0.2}{\second}$.
              因此$u_T \approx \sqrt{(\qty{0.1}{\second})^2 + (\qty{0.2}{\second})^2} \approx \qty{0.22}{\second}$.
    \end{itemize}
\end{rhoenv}

\section{测量周期数}

为了使实验精度达到$\frac{u_g}{g} \leqslant 1 \qty{1}{\percent}$,我们需要确定重复测量周期数$N$.
当摆长$l = \qty{70}{\centi\metre}$时,取$g = \qty{9.8}{\metre\per\second\squared}$,得单次周期$T = \qty{1.68}{\second}$.若仅测量单次周期,相对不确定度约为$\qty{23.8}{\percent}$.

考虑重复测量$N$个周期,此时不确定度变为
\begin{equation}
    \frac{u_g}{g} = \sqrt{\left( \frac{u_l}{l} \right)^2 + \left( \frac{2u_T}{NT} \right)^2} \approx \sqrt{\left( 0.00286 \right)^2 + \left( \frac{0.262}{N} \right)^2}
\end{equation}

令$\frac{u_g}{g} \leqslant 1 \qty{1}{\percent}$,解得$N \geqslant 25$,即至少要测量 25 个周期.

\end{document}